\documentclass[FGBook.tex]{subfiles}

\begin{document}

\poezie{Máj}{Karel Hynek Mácha}

\noindent\begin{wrapfigure}{r}{0.35\textwidth}\footnotesize
\subwection{Ukázka z díla:}
\setlength{\parindent}{3pt}
\begin{center}
\noindent
Vyrazil jsem okenici, když tu s velkou motanicí \\
vstoupil starodávný havran z dob, jež jsou tak záslužné; \\
bez poklony, bez váhání, vznešeně jak pán či paní \\
usadil se znenadání v póze velmi výhružné \\
na poprsí Pallady - a v póze velmi výhružné \\
si sedl jen a víc už ne. 
\end{center}
\end{wrapfigure}

\wection{LITERÁRNÍ TEORIE}

\subwection{Literární druh a žánr:}
\noindent 
Lyricko-epická báseň, spíše balada.

% \subwection{Literární směr:}

% \subwection{Jazyk a slovní zásoba:}

% \subwection{Figury:}

% \subwection{Tropy:}

\subwection{Stylistická charakteristika textu:}
\noindent 
Jedná se o krátkou barvitou báseň, kdy autor užívá popisů pro navození atmosféry temného pokoje.
Je zde kontrast mezi klidným pokojem a bouří za okny.
I přes krátkost básně se autorovi daří bravurně popsat atmosféru i psychické rozpoložení hlavní postavy,
která je promítnutím samotného autora.

\subwection{Postavy:} 
\noindent
\textschl{Vypravěč --} Autor sám, truchlí nad smrtí Lenory, své milované,
stejně jako autor truchlil nad smrtí své ženy.\\
\textschl{Havran --} Tajemný pták, co přiletí z bouřlivé noci a opakuje jedno slovo,
které ho klidně mohl naučit majitel, \enquote{nevermore}
(v překladu jde většinou o dvou či tří slovní frázi, v mé verzi \enquote{nikdy více}).
Jak autor propadá ve svých myšlenkách hloubš a hloubš, připisuje havranovým slovům větší a větší smysl.\\
\textschl{Lenora --} Mrtvá milenka, není zde o ní moc řečeno, než to, že nad ní hlavní postava truchlí.

\subwection{Děj:}
{\setlength{\parindent}{0pt}\setlength{\parskip}{0.5em}
\noindent Ocitáme se v pokoji učence, starého muže, kterého trápí ztráta milenky
a snaží se zapomenout s nosem v knížce. Zatímco je muž uvnitř svého malého kruhu míru,
tak za okny zuří bouřka.\\
Ťuk, ťuk, ťuk. Muže z dlení vytrhne zaťukání, poleká se, racionálně se snaží přesvědčit,
že o nic nejde, jde otevřít dveře a tam nikdo.\\
Ťuk, ťuk, ťuk. Zlověstné zaťukání bez původu. Muž otevírá okenici,
čímž je jeho klid narušen už ničím nezjemněnou bouří.
Z tohoto nehezkého prostředí do pokoje vlétá Havran.
(V originále Krkavec, ale překladatelé se rozhodli, že nevíme, co krkavec je
a že havran je u nás více v povědomí jako stvoření Satanovo, čehož autor se snaží docílit.
Nevíme, zda byl poslán z druhé strany, zda jenom blouzníme ze zmařené lásky společně s autorem,
či zda k nám opravdu přišel posel se zprávou. Havran, i přes veškerou symboliku, může být jen ptákem,
který se chtěl jít schovat do teplého příbytku před bouří. Také ale mohl muži přinést důležitou zprávu.)
Havran se usadí na poprsí Pallady (socha Pallas Athény) což muži připadá ironické a i se tomu zasměje.
Další várka symbolismu.\\
Muž se ptá na různé otázky, kdy vždy dostane odpověď zápornou a konečnou, \enquote{nevermore}.
Ptá se dál, kdy vždy dostává stejnou odpověď a propadá čím dál tím hloubš do sebe trýzně.
}

\subwection{Kompozice, prostor a čas:}
\noindent 
Báseň o krátkém časovém období, psána chronologicky. 18 slok a pouze 108 veršů.
Jde-li skutečně o prakticky auto-biografii, vypsání z depresí básníka, tak by se báseň odehrávala
v devatenáctém století.

% \subwection{Význam sdělení:}

\wection{LITERÁRNÍ HISTORIE}

% \subwection{Politická situace:}

% \subwection{Základní principy fungování společnosti:}

% \subwection{Kontext dalších druhů umění:}

% \subwection{Kontext literárního vývoje:}

\subwection{Edgar Allan Poe {\ssmall -- život autora:}}
\noindent 
Autor vedl celkem bídný život.
Otec alkoholik se na ně vykašlal, matka zemřela, když byl mladý.
Byl umístěn do dětstkého domova, odkud si ho převzal bohatý obchodník a umožnil mu studium.
Pár let žili spolu i v Anglii, která se mu mnohokrát stala inspirací pro svojí práci.
Díky své povaze a možná i vzpomínkách byl však vysoce nevyrovnaný.
Propadal hráčství a závislostem, i když i za života sklidil slávu, tak nikdy nebyl přímo bohatý.
Nejspíš i toto mu pomohlo napsat jeho skvosty.\\
Edgar Allan Poe byl i kritizován za své manželství se svojí sestřenicí.
V době svatby jemu bylo sedmadvacet a jí třináct.
Dle určitých zdrojů se jejich vztah podobal spíše sourozeneckému než manželskému.
Virginia Poe zemřela mladá na nemoc.

% \subwection{Vlivy na dané dílo:}

\subwection{Další autorova tvorba:}
\noindent Mezi další autorovu tvorbu patří \textsc{Rukopis nalezený v láhvi},
\textsc{Filozofie básnické skladby} a \textsc{Příběhy Arthura Gordona Pyma z Nantucketu}.

% \subwection{Jak dílo inspirovalo další vývoj literatury:}

% \wection{LITERÁRNÍ KRITIKA}

% \subwection{Dobové vnímání díla a jeho proměny:}

% \subwection{Aktuálnost tématu a zpracování tématu:}


\wection{OSOBNÍ NÁZOR}
\noindent 
Havran je bravurně zpracovanou básní, která zvládá dokonale obrazutvornost. I týdny po přečtení si jsem schopen vybavit pochmurnost básníkova pokoje a jak jsem soucítil s jeho trápením, které jsem ani nezažil.

\newpage

\end{document}